\chapter{Plastics, Rubber, and Fibers}

\begin{kaichang}
Bring three things from around the house: an empty water bottle, a rubber band, and a garment labeled polyester or polyester fiber. A supermarket plastic bag can substitute for the garment.

Guess two things. First, one is hard, one soft, and one drawable into fine threads, yet chemists call them one class of materials. Why? Second, why does a plastic water bottle resist shattering and decay in soil, while a rubber band left in sunlight for years turns sticky, cracks, and snaps?

Write your guesses. By the end, one extremely long thread will explain all three, including why polyester clothing labels often share the abbreviation PET.
\end{kaichang}

\section{Light, Tough, and Reluctant to Decay}

Start only with facts you can see and touch.

A 500-milliliter glass bottle weighs about three or four hundred grams, while a comparable plastic bottle weighs roughly twenty and does not shatter. Transparent, bendable, and twistable, it can hold water for a year without dissolving or molding. Bury it and dig it up decades later, and its shape remains: both virtue and problem.

A rubber band is the opposite: soft, stretchable to three or four times its length, and snapping back on release. Yet a summer forgotten on a windowsill leaves it sticky, brittle, and easily broken.

Polyester clothing is light, thin, quick-drying, and nearly wrinkle-free. Near a barbecue, however, a spark does not simply scorch it: it melts a small hard-edged hole.

Tough but persistent, elastic but aging, melting rather than burning: three apparently unrelated observations share one microscopic structure, \textbf{extremely long molecular chains}.

\section{Long Chains of Small Molecules}

Chapter 36 gave carbon at most four hands, covalent bonds, allowing chains, rings, and networks. Earlier organic molecules, ethanol, acetic acid, aspirin, held only dozens of atoms. Now imagine one carbon joining another and another, repeated thousands of times.

Start with tiny ethene, two carbons and four hydrogens, with a double bond between carbons. Chemical plants join thousands of ethene molecules, opening doubles into a long chain of thousands of carbons. Ethene is the \textbf{monomer}, a single small molecule; the chain is a \textbf{polymer}, a joined multitude. Its familiar commercial name is polyethylene, the material of bags and milk jugs.

A polyethylene molecule has thousands of carbons but \textbf{no} chemical bonds to neighboring molecules. Between chains lie Chapter 22's weak but numerous intermolecular forces. This unlocks plastic behavior: imagine cooked noodles. Each noodle is internally strong, carbon--carbon covalent bonds costly to break (Chapter 19), while noodles slide and entangle. Pull, and they first straighten, then slide apart, deforming the material. Heat increases molecular motion and sliding, so it softens and melts. Plasticity means chains sliding past each other, not chains themselves breaking.

Resistance to decay makes sense too. The carbon--carbon backbone offers no interface microbial enzymes recognize, and water cannot attack it. Degradation is not thermodynamically forbidden but immeasurably slow: thermodynamics sets direction, kinetics speed, Chapter 29's rule.

\section{Two Ways to Join: Addition and Condensation}

Chemists mainly join small molecules by two methods, differing in their interfaces.

\textbf{Addition polymerization forms a relay chain}. Double-bonded monomers, ethene, vinyl chloride, styrene, open their doubles and join hand in hand, leaving single bonds along the chain. All starting atoms enter it, with no byproducts. An initiator provides the first push. Older high-temperature, high-pressure radical processes produce highly branched low-density polyethylene, LDPE, soft and transparent like a bag. Later Ziegler and Natta developed special catalysts: Chapter 30 says catalysts change pathways, not direction or equilibrium. Lower-temperature, lower-pressure orderly straight chains yield high-density polyethylene, HDPE, hard and milky white, used in milk jugs and drainpipes. Same monomer, different start-up method, entirely different material, a contrast we will reuse.

\textbf{Condensation polymerization joins paired interfaces}. Each monomer has two Chapter 37 functional groups. Meeting groups form a bond and release a small molecule, usually water. Chapter 40 joined acid and alcohol into ester plus water; now apply it to molecules functional at both ends. Terephthalic acid, two carboxyls, and ethylene glycol, two hydroxyls, join by an ester and release water, leaving ends free for further joining. The growing polyester is PET, made into drink bottles or drawn into polyester fiber. Substitute adipic acid and amino-ended hexamethylenediamine, and amide links produce nylon.

Remember the accounting: addition retains every entering atom in the chain; condensation pays a water fee at each joining.

\begin{qiaoliang}
How long is a polymer chain? Take polyethylene with molar mass about \ud{140000}{g\;mol^{-1}}. Each ethene unit contributes \ud{28}{g\;mol^{-1}}, giving $140000\div 28\approx 5000$ units, about ten thousand carbon--carbon single bonds. At \ud{0.154}{nm} each, a straightened chain spans $10000\times 0.154\approx 1540$ nanometers, about 1.5 micrometers. Its zigzag makes actual extended length a little over one micrometer, the same order. Normally it coils randomly into only tens of nanometers diameter. Thus \textbf{a straightened chain is tens of times longer than its usual occupied space}. A small pinch of plastic holds over $10^{18}$ intertwined chains: that noodle bowl's real density.
\end{qiaoliang}

\section{Four Controls for Material Properties}

Chemists have four controls turning the same ingredients into very different materials. All begin with strong covalent bonds inside chains and weak, sliding interactions between them.

\begin{itemize}[itemsep=2pt]
\item \textbf{Chain length}: longer chains entangle more firmly, increasing toughness and resistance to melt flow. Mini-chains of twenty or thirty carbons make crumbly candle wax; thousands make bags; tens of thousands in ultrahigh-molecular-weight polyethylene make body armor and artificial joints.
\item \textbf{Branching}: highly branched LDPE resembles twig-covered branches that pack poorly, giving weak interchain attraction, softness, low melting point, and transparency. Nearly unbranched HDPE resembles aligned chopsticks, packed tightly into hard, strong, milky-white material.
\item \textbf{Crystallinity}: orderly chain segments form small crystalline regions, recalling Chapter 24. Like nails, these pin surrounding chains. More crystallinity generally makes materials harder, stronger, and more heat-resistant, but scattering at crystal boundaries reduces transparency. PET bottles are clear and tough because factories hold crystallinity appropriately low.
\item \textbf{Cross-linking}: real chemical bonds bridge chains. A few let segments slide and straighten while tethering chains so they spring back: rubber. Many weld them into an immobile three-dimensional network: Bakelite sockets and epoxy, which char before melting.
\end{itemize}

The last control divides \textbf{thermoplastics}, with no or very few cross-links, repeatedly softening hot and hardening cold, from \textbf{thermosets}, cross-linked networks that decompose and char rather than melt. Bottles and bags are thermoplastic; tires and socket housings thermoset. This distinction will decide which can be recycled and which only landfilled.

Rubber deserves special attention. Natural rubber comes from tree latex, with isoprene monomers and naturally coiled chains, leaving raw rubber floppy, sticky, and easily torn. In 1839, American Goodyear accidentally dropped sulfur-mixed rubber onto a hot stove and found it dry and elastic: sulfur built a few bridges, vulcanization. Few enough for segments to straighten, enough to stop chains escaping. Modern tires use this principle with far more refined formulations.

Another plastic unlike plastics is silicone, in nipples, spatulas, baking mats, and phone cases. Its backbone alternates silicon and oxygen, \chem{Si{-}O{-}Si{-}O}, rather than carbon, with two small organic groups on silicon. Heat-resistant, flexible silicon--oxygen bonds let it go in oven or refrigerator, reminding us polymer backbones need not be carbon.

One final ethene variation replaces all chain hydrogens with fluorine, giving polytetrafluoroethylene, PTFE or Teflon, the glossy black nonstick-pan coating. Strong carbon--fluorine bonds and a close-fitting fluorine armor prevent other molecules gripping or attacking the chain: food does not stick, and strong acids and bases cannot harm it. Yet it is not invincible: an empty pan heated beyond two or three hundred degrees starts chain decomposition. Proper use means no empty heating or metal-spatula scraping, not anything goes.

Why is rubber soft at room temperature while a PET bottle shrinks and deforms with boiling water? Each chain has a freezing switch, its glass-transition temperature, related to Chapter 24's glass. Below it segments cannot move and material is hard and brittle; above, they twist and soften. PET's switch is around $70\,^{\circ}\mathrm{C}$, crossed by boiling water. Rubber's is around minus sixty or seventy, leaving it soft all winter.

\section{Now, the Symbols}

After so many chains, write equations.

Addition polymerization opens $n$ ethene double bonds and joins a chain:

\[n\,\chem{CH_2{=}CH_2} \rightarrow \chem{{[CH_2{-}CH_2]}_{n}}\]

Brackets contain the \textbf{repeating unit}; subscript $n$ counts repeats. Replace one ethene hydrogen with chlorine to get vinyl chloride, then polymerize to PVC for pipes and floor coverings:

\[\chem{{[CH_2{-}CHCl]}_{n}}\]

Rigid PVC makes pipes; abundant plasticizer, small molecules entering between chains to ease sliding, makes flexible wire insulation and table mats.

Keep condensation as word accounts: terephthalic acid $+$ ethylene glycol $\rightarrow$ PET $+$ water; adipic acid $+$ hexamethylenediamine $\rightarrow$ nylon $+$ water. Each ester or amide link releases one \chem{H_2O}.

Natural rubber's repeating unit derives from isoprene, \chem{CH_2{=}C(CH_3){-}CH{=}CH_2}. One double bond remains per unit after polymerization, vulnerable to oxygen and ozone, making sun-aged bands sticky and broken. Tire carbon black and antioxidants protect these weak points.

One fact must be explicit: chain $n$ values differ within a bottle, perhaps three thousand here and five thousand there. Polymers have no single molecular mass, but an average and distribution. Polymeric materials are inherently mixtures, unlike Volume I's pure water or sodium chloride.

\begin{zhengju}
Thousands of atoms per molecule seems obvious today, but was fiercely debated in the 1920s. Mainstream opinion called rubber, starch, and cellulose colloids: small molecules loosely clustered by mysterious residual valence forces, not genuine macromolecules. X-ray unit cells were small, seemingly supporting small molecules.

In 1920, German chemist Staudinger proposed the opposite: rubber is isoprene joined by ordinary covalent bonds into long chains, with no mysterious forces. He designed a decisive experiment. If double-bond residual forces held small molecules together, hydrogenating all doubles should destroy the glue and break clusters apart. If covalent long chains were real, removing doubles would leave chains long. Hydrogenated rubber retained macromolecular behavior, refuting the imagined glue.

Quantitative evidence followed. Staudinger found more viscous solutions of a given polymer indicated longer chains, making viscosity a molecular-chain scale. Osmotic pressure and later ultracentrifugation gave molecular masses from thousands to hundreds of thousands, two orders above what anyone wished to believe. Debate continued into the 1930s until DuPont's Carothers synthesized well-defined nylon step by condensation step in 1935. Humans could build macromolecules piece by piece, finally retiring the colloid school. Nylon stockings launched in 1940, selling hundreds of thousands of pairs on day one; Staudinger received the Chemistry Nobel in 1953.

The point is not who was clever. Both explanations fitted tough, sticky material; deciding required \textbf{an experiment making them predict different outcomes}. Hydrogenation provided that fork.
\end{zhengju}

\textbf{Translator safety note:} Do not put a stretched rubber band against your lip or near your eyes; skip that contact step. Natural rubber may contain latex allergens. Keep the activity away from the face and avoid latex if allergic. See \href{https://www.cdc.gov/niosh/docs/2012-119/default.html}{NIOSH latex-allergy guidance}.

\begin{shiyan}
\textbf{Stretch and Watch Chains Line Up}

Materials: a clean polyethylene shopping bag, confirmed PE, and a rubber band.

Procedure: cut a bag strip about two centimeters wide and fifteen long. Hold its ends and pull slowly and steadily, watching the middle. Stop after some extension and observe changes. Then stretch the band slowly and lightly touch it to your temperature-sensitive upper lip for two or three seconds before releasing.

What to observe: the bag strip suddenly develops a \textbf{thinner, whiter} neck, spreading toward the ends with further stretching. This necking straightens and rearranges disordered chains, leaving aligned regions unable to contract back. Rubber differs completely: stretching feels slightly warm, and release restores length.

Safety: hands only, no heating or flames. \textbf{Never} burn plastic to observe it: fumes can irritate or poison. For discarded-plastic handling, inspect recycling categories or watch a teacher's demonstration video. Put used strips in a recyclables bin.
\end{shiyan}

\section{Where Does This Model Fail?}

Strong intrachain bonds, weak interchain forces, four controls: this is the backbone of polymer understanding, but real products are more complex.

First, almost no commercial plastic is pure polymer. PVC has plasticizers, tires carbon black, bottles processing stabilizers, colored plastics pigments. Formulations are engineering, often publicly listing only main ingredients. PET alone does not tell everything.

Second, molecular mass is distributed and crystallinity averaged, with crystalline and amorphous regions interwoven. Models predict trends, such as branching softening, not a specific bottle's exact strength.

Third, degradation is deceptive. Sunlight can embrittle and fragment a bag, shortening chains and making microparticles without returning carbon to natural cycles. Invisible is not gone. Most labeled degradable plastics need hot, humid industrial composting to truly break down; in grass or oceans they may persist like ordinary plastics for years. Environmental accounts require raw materials, production, use, and disposal together: life-cycle assessment, coming shortly.

\begin{wuqu}
``Polymers are all synthetic plastics; nature contains none.'' Wildly wrong, missing half this book. Starch and cellulose are glucose polymers, whose different digestibility Chapter 46 explains; proteins join twenty amino acids (Chapter 48); DNA joins four nucleotides into extraordinarily long chains. Humans have made plastics for a little over a century; nature has joined monomers for more than three billion years. Chemists truly invented only the use of petroleum-derived monomers.
\end{wuqu}

\section{Back to the Bottle}

Now fulfill the opening promise.

The water bottle is PET from acid and glycol joining while paying water fees, about one or two hundred units per chain. It is light because chains have low density and moderate packing; shatter-resistant because tens of thousands of entangled chains slide to absorb impact rather than propagate glass-like cracks; clear because crystallinity is kept low; strong because bottle blowing stretches and aligns chains in two directions, like the bag's necking. It resists decay not from bacterial reluctance but absent recognized interfaces on its carbon framework and nearly imperceptibly slow room-temperature ester hydrolysis.

The rubber band is polyisoprene plus a few sulfur bridges. Coiled chains straighten, bridges stop escape and permit recoil, and oxygen, ozone, and ultraviolet attack remaining doubles, puncturing the network and causing sticky breakage.

Polyester clothing is the bottle's PET melted through fine holes and drawn into threads. Polyester, polyester fiber, and PET thus name the same chain. Few water-gripping sites make it quick-drying; sparks melt holes because thermoplastic chains soften and flow hot.

Three objects, one logic: covalent bonds inside chains determine strength; interchain forces and cross-links determine movement, which nearly determines every feel under your fingers.

\section{Digging Deeper: A Bottle's Afterlife}

Look beneath a bottle for a number inside a triangle, plastic's recycling identity card:

\begin{table}[htbp]
\centering
\begin{tabular}{clll}
\toprule
Code & Material & Common examples & Recycling destination \\
\midrule
1 & PET (polyester) & \shortstack[l]{Drink bottles,\\cooking-oil bottles} & \shortstack[l]{Mature recycling;\\can become fiber} \\
2 & \shortstack[l]{HDPE (high-density\\polyethylene)} & \shortstack[l]{Milk jugs,\\shampoo bottles} & \shortstack[l]{Fairly mature\\recycling} \\
5 & PP (polypropylene) & \shortstack[l]{Takeout containers,\\bottle caps} & Recycled in some areas \\
6 & PS (polystyrene) & \shortstack[l]{Foam food boxes,\\disposable cup lids} & Difficult to recycle \\
7 & \shortstack[l]{Other (including\\PC, acrylic, etc.)} & \shortstack[l]{Water jugs, lenses,\\mixed materials} & Generally not recycled \\
\bottomrule
\end{tabular}
\end{table}

Behind it is an engineering problem: recycling must separate polymers with different melting points and incompatible chains, or regenerated material performs poorly. Thermoplastics at least remelt; thermoset tires and circuit boards cannot clear even that step. Much plastic therefore reaches landfill, incineration, or the environment, where sun and waves fragment it into microplastics under five millimeters, found in deep seas, soil, and even blood. Health effects remain under study: high laboratory doses can harm cells and animals, but actual ingested doses and consequences lack sufficiently strong evidence. This warrants concern without definitive conclusions; apply the dose principle to headlines.

Why not abandon plastics for paper and cotton? Life-cycle accounts matter: paper bags are heavier and use more water and trees; a cotton bag's carbon bill far exceeds a plastic one's, with some studies requiring tens to hundreds of reuses to break even climatically. The answer is not banning one material, but using every chain long enough and recycling cleanly enough.

One mystery remains. Humans use two chain-making methods, while nature does more: one glucose monomer produces digestible starch and indigestible cellulose; twenty amino acids join in precise sequences into self-folding machines. How can chains differ as much as fuel and machine? Starting Chapter 45, enter nature's polymer factory. Monomers, condensation, interchain forces, and cross-links return, but its ledger follows life's own algorithms.

\begin{tiaozhan}
Why does stretched rubber warm? Counterintuitively, elasticity mainly comes not from stretched bonds recoiling but \textbf{entropy}, Chapter 28. Coiled chains have enormous microscopic arrangement counts and high entropy; straightened chains have far fewer and low entropy. Released chains randomly jostle toward high-entropy coils, appearing as contraction. Forced ordering on stretching releases thermal-motion energy as heat, creating warmth. Conversely, heating a stretched band contracts it instead of expanding it like most solids. This predicts a test: suspend rubber at tens of degrees below zero, and crossing its freezing temperature makes it brittle glass. Thermodynamic statistics hides in your pencil case.
\end{tiaozhan}

\section*{Try It Yourself}

\begin{sikaoti}
\item Draw material-flow diagrams for $n$ ethene molecules becoming polyethylene, marking opened doubles and whether byproducts occur, and acid plus glycol condensing into PET, marking where water leaves. Compare atom accounts.
\item Using only chain length, branching, cross-linking, and crystallinity, explain milky rigid HDPE milk jugs versus transparent soft LDPE bags despite identical monomers.
\item A PVC chain averages a thousand repeating units, each \ud{62.5}{g\;mol^{-1}}. Estimate molar mass and extended length at 0.25 nanometers per unit; compare its coiled size of tens of nanometers.
\item Bottles remelt into fiber while tires can only be shredded for roads. Sketch chains as lines and cross-links as short bridges to explain why thermosets cannot melt.
\item Starch and cellulose both join glucose, but humans digest only starch. Using enzyme-interface recognition, predict the key to digestibility in a natural polymer. Chapter 46 answers.
\end{sikaoti}
